\chapter{统一时间与离散事件}
\section{逻辑时钟}
状态保存逻辑时间\code{now}。正常原生操作使用\code{Date.now()+timeOffset}；休整增加偏移并向未来推进。目标向下取整并钳制为不早于当前时间，系统时钟回拨不会让游戏状态直接倒退。

\code{advance()}克隆状态、调用\code{advanceInPlace()}并增加revision。订单、采集和刷新不依赖帧数，而是跳到最近事件边界，在边界间累计自产和采集。

\section{排序与分段累计}
\begin{small}\noindent\begin{tabularx}{\linewidth}{@{}P{27mm}P{19mm}Y@{}}
\toprule 事件 & 优先级 & 同时刻行为\\\midrule
订单交付 & 10 & 按seq排序；逐辆交付或在免费整批边界交付剩余量\\
采集完成 & 20 & 处理尾量并转返程\\
去程/返程到达 & 30 & 战斗、收据、返还车辆或入库\\
世界补给 & 独立比较 & 取地点下一刷新时刻；与个人事件相等时个人事件先处理\\\bottomrule
\end{tabularx}\end{small}

每轮取个人事件与世界事件的最小时间。超过目标时，增长剩余时段并结束；否则先增长至边界，再处理事件。建筑或科技完成后的时段使用新产率和容量。因此不能先把整段离线时间都按最终等级计算。

自产保存\code{remainders}，采集保存\code{March.remainder}。以BigInt计算“每小时速率乘毫秒加余数”，除以3600000得到整数收益；余数用于下一段。

\section{逐辆交付与免费边界}
多辆订单的\code{duration}是单辆工时，\code{dueAt}是下一辆完成时间。总剩余为该辆剩余加后续车辆工时。VIP免费分钟数作用于整单剩余边界，不是给每辆重复减一次。

活动订单结束后，同工位最早等待单在当前时刻开工，马上再检查免费完成条件。目标等于当前时刻也会排空立即事件，短订单因此可能提交即完成。

\section{休整与规模限制}
休整输入为1至720整数小时，保存在Godot偏好中。单次推进跨度和累计休整偏移有防御性上限；事件循环超过100000次中止。它们不是无限长挂机的性能承诺。

候选每轮重新收集排序，世界刷新检查100个地点。工位数有限，适合当前单机规模；并未实现持久优先队列。很长推进和大存档可能增加同步计算耗时。
\src{src/core/engine.ts}{grow；advanceInPlace；finishJob；resolveArrival}
